\begin{table}[H]
\caption{Características de la cohorte según diagnóstico angiográfico (n = 297). Continuas: media (DE) o mediana [RIC] según Shapiro-Wilk; cualitativas: n (\%).}
\label{tab:tabla1}
\begin{tabular}{lrrrrrrr}
\toprule
Variable & Categoría & Total & Sin EC & Con EC & Prueba & p & Efecto \\
\midrule
n &  & 297 & 160 & 137 &  &  &  \\
Edad (años) &  & 56,0 [48,0–61,0] & 52,0 [44,8–59,0] & 58,0 [53,0–62,0] & U de Mann-Whitney & < 0,001 & r\_rb = 0,28 \\
Presión arterial sistólica en reposo (mmHg) &  & 130,0 [120,0–140,0] & 130,0 [120,0–140,0] & 130,0 [120,0–145,0] & U de Mann-Whitney & 0,023 & r\_rb = 0,15 \\
Colesterol sérico total (mg/dL) &  & 243,0 [211,0–276,0] & 235,5 [208,8–268,2] & 253,0 [218,0–284,0] & U de Mann-Whitney & 0,047 & r\_rb = 0,13 \\
Frecuencia cardíaca máxima alcanzada (lpm) &  & 153,0 [133,0–166,0] & 161,0 [149,0–172,0] & 142,0 [125,0–157,0] & U de Mann-Whitney & < 0,001 & r\_rb = -0,50 \\
Depresión del ST inducida por ejercicio (mm) &  & 0,8 [0,0–1,6] & 0,2 [0,0–1,1] & 1,4 [0,6–2,5] & U de Mann-Whitney & < 0,001 & r\_rb = 0,47 \\
Vasos mayores coloreados por fluoroscopia (n.º (0-3)) &  & 0,0 [0,0–1,0] & 0,0 [0,0–0,0] & 1,0 [0,0–2,0] & U de Mann-Whitney & < 0,001 & r\_rb = 0,50 \\
Sexo masculino & Mujer & 96 (32,3) & 71 (44,4) & 25 (18,2) & Chi² de Pearson & < 0,001 & V = 0,28 \\
 & Hombre & 201 (67,7) & 89 (55,6) & 112 (81,8) &  &  &  \\
Glucemia en ayunas > 120 mg/dL & No & 254 (85,5) & 137 (85,6) & 117 (85,4) & Chi² de Pearson & 0,956 & V = 0,00 \\
 & Sí & 43 (14,5) & 23 (14,4) & 20 (14,6) &  &  &  \\
Angina inducida por ejercicio & No & 200 (67,3) & 137 (85,6) & 63 (46,0) & Chi² de Pearson & < 0,001 & V = 0,42 \\
 & Sí & 97 (32,7) & 23 (14,4) & 74 (54,0) &  &  &  \\
Tipo de dolor torácico & Angina típica & 23 (7,7) & 16 (10,0) & 7 (5,1) & Chi² de Pearson & < 0,001 & V = 0,51 \\
 & Angina atípica & 49 (16,5) & 40 (25,0) & 9 (6,6) &  &  &  \\
 & Dolor no anginoso & 83 (27,9) & 65 (40,6) & 18 (13,1) &  &  &  \\
 & Asintomático & 142 (47,8) & 39 (24,4) & 103 (75,2) &  &  &  \\
ECG en reposo & Normal & 147 (49,5) & 92 (57,5) & 55 (40,1) & Chi² (Monte Carlo) & 0,005 & V = 0,18 \\
 & Anomalía ST-T & 4 (1,3) & 1 (0,6) & 3 (2,2) &  &  &  \\
 & Hipertrofia VI & 146 (49,2) & 67 (41,9) & 79 (57,7) &  &  &  \\
Pendiente del ST en el pico de ejercicio & Ascendente & 139 (46,8) & 103 (64,4) & 36 (26,3) & Chi² de Pearson & < 0,001 & V = 0,38 \\
 & Plana & 137 (46,1) & 48 (30,0) & 89 (65,0) &  &  &  \\
 & Descendente & 21 (7,1) & 9 (5,6) & 12 (8,8) &  &  &  \\
Gammagrafía con talio & Normal & 164 (55,2) & 127 (79,4) & 37 (27,0) & Chi² de Pearson & < 0,001 & V = 0,53 \\
 & Defecto fijo & 18 (6,1) & 6 (3,8) & 12 (8,8) &  &  &  \\
 & Defecto reversible & 115 (38,7) & 27 (16,9) & 88 (64,2) &  &  &  \\
\bottomrule
\end{tabular}
\end{table}
